\cleardoublepage
\markboth{CAPÍTULO IV}{CAPÍTULO IV}
\oldpage[IV-1]{222}
\thispagestyle{plain}

\begin{center}
{\Large CAPÍTULO IV}

\vspace{1.5cm}
{\LARGE\bfseries ESTUDIO LOCAL DE LOS ESQUEMAS\\[0.3em]
Y DE LOS MORFISMOS DE ESQUEMAS}

\vspace{1.5cm}
{\large\bfseries Sumario}
\end{center}

\medskip
\begin{tabular}{@{}r p{0.84\textwidth}}
§ 1. & Condiciones de finitud relativas. Conjuntos construibles en los preesquemas.\\
§ 2. & Cambio de base y planitud.\\
§ 3. & Ciclos primos asociados y descomposiciones primarias.\\
§ 4. & Cambio del cuerpo de base en los preesquemas.\\
§ 5. & Dimensión y profundidad en los preesquemas.\\
§ 6. & Morfismos planos de preesquemas localmente noetherianos.\\
§ 7. & Aplicación a las relaciones entre un anillo local noetheriano y su completado.
Anillos excelentes.\\
§ 8. & Límites proyectivos de preesquemas.\\
§ 9. & Propiedades construibles.\\
§ 10. & Preesquemas de Jacobson.\\
§ 11.\footnotemark[1] & Propiedades topológicas de los morfismos planos de presentación finita.
Criterios locales de planitud.\\
§ 12. & Estudio de las fibras de los morfismos planos de presentación finita.\\
§ 13. & Morfismos equidimensionales.\\
§ 14. & Morfismos universalmente abiertos.\\
§ 15. & Estudio de las fibras de un morfismo universalmente abierto.\\
§ 16. & Invariantes diferenciales. Morfismos diferencialmente suaves.\\
§ 17. & Morfismos suaves, no ramificados, étales.\\
§ 18. & Complementos sobre los morfismos étales. Anillos locales henselianos.\\
§ 19. & Inmersiones regulares y transversalmente regulares.\\
§ 20. & Secciones hiperplanas; proyecciones genéricas.\\
§ 21. & Prolongaciones infinitesimales.\\
\end{tabular}
\footnotetext[1]{El orden y el contenido de los §§ 11 a 21 se dan únicamente a título
indicativo y podrán modificarse antes de su publicación.}

\clearpage
\oldpage[IV-1]{223}

Los temas tratados en este Capítulo requieren las observaciones siguientes:

\medskip
\noindent
\emph{a)} Su carácter común es referirse a propiedades
\emph{locales} de preesquemas o de morfismos, i.e. consideradas en un punto,
o en los puntos de una fibra, o en una vecindad (no especificada) de un punto o
de una fibra. Estas propiedades son generalmente de naturaleza \emph{topológica},
\emph{diferencial} o \emph{dimensional} (i.e. en las que intervienen las
nociones de \emph{dimensión} y de \emph{profundidad}) y están ligadas a las
propiedades de los \emph{anillos locales} en los puntos considerados. Un problema típico
es relacionar, para un morfismo dado $f:X\to Y$ y un punto
$x\in X$, las propiedades de $X$ en $x$ con las de $Y$ en $y=f(x)$ y de
la fibra $X_y=f^{-1}(y)$ en $x$. Otro es determinar la naturaleza
topológica (por ejemplo la construibilidad, o el hecho de ser abierto o
cerrado) del conjunto de los puntos $x\in X$ en los que $X$ posee cierta
propiedad o para los cuales la fibra $X_{f(x)}$ que pasa por $x$ posee una
cierta propiedad en $x$. Asimismo, se estudiará la naturaleza topológica
del conjunto de los puntos $y\in Y$ tales que $X$ tenga cierta propiedad en
todos los puntos de la fibra $X_y$, o tales que esta fibra tenga cierta
propiedad.

\medskip
\noindent
\emph{b)} Las nociones más importantes para los capítulos siguientes son
la de \emph{morfismo plano de presentación finita} y las de
\emph{morfismo suave} y de \emph{morfismo étale}, que son casos
particulares de aquella. Su estudio detallado (junto con el de cuestiones conexas)
comienza propiamente en el §~11.

\medskip
\noindent
\emph{c)} Los §§~1 a 10 pueden considerarse de naturaleza preliminar
y desarrollan tres tipos de técnicas, utilizadas no solo en los
otros parágrafos del capítulo sino también, por supuesto, en los capítulos
siguientes:

\medskip
\noindent
\emph{c} 1) En los §§~1 a 4 se consideran diversos aspectos de la noción de
\emph{cambio de base}, sobre todo en relación con condiciones de
\emph{finitud} o de \emph{planitud}; se inicia allí la técnica de
\emph{descenso} en sus aspectos más elementales (las cuestiones
de «efectividad» ligadas a esta técnica se estudiarán en el cap.~V).

\medskip
\noindent
\emph{c} 2) Los §§~5 a 7 se centran en lo que puede llamarse
técnicas \emph{noetherianas}, suponiéndose siempre que los preesquemas considerados
son localmente noetherianos, mientras que, por el contrario, generalmente no se
impone ninguna condición de finitud a los \emph{morfismos}; esto se debe
esencialmente a que las nociones de dimensión y de profundidad apenas son
manejables fuera de los anillos locales noetherianos. Recordemos que el §~7
constituye una teoría «fina» de los anillos locales noetherianos, bastante poco
utilizada en el resto del capítulo.

\medskip
\noindent
\emph{c} 3) Los §§~8 a 10 proporcionan, entre otras cosas, el medio de \emph{eliminar las
hipótesis noetherianas} sobre los preesquemas estudiados, sustituyéndolas por
hipótesis adecuadas de \emph{finitud} (la «presentación finita») sobre los
\emph{morfismos} considerados: la ventaja de esta sustitución es que estas
últimas hipótesis son \emph{estables por cambio de base}, lo que no
ocurre con las hipótesis noetherianas sobre los preesquemas. La técnica
que permite este paso se basa, por una parte, en la utilización de la noción de
\emph{límite proyectivo} de preesquemas, gracias a la cual se puede reducir una
cuestión a la misma cuestión bajo hipótesis \emph{noetherianas}; por otra
parte, en el empleo sistemático de los \emph{conjuntos construibles}, que tienen
el doble interés de conservarse por imagen inversa (por un morfismo
\oldpage[IV-1]{224}
cualquiera) y por imagen directa (por un morfismo de presentación finita), y
de tener propiedades topológicas manejables en los preesquemas localmente
noetherianos. Las mismas técnicas permiten incluso a menudo reducirse a
anillos noetherianos más particulares, por ejemplo las
$\mathbf Z$-\emph{álgebras de tipo finito}, y es ahí donde las propiedades de los
anillos «excelentes» (estudiadas en el §~7) intervienen de manera decisiva.
Por lo demás, independientemente de la cuestión de eliminar las hipótesis
noetherianas, las técnicas de los §§~8 a 10, de naturaleza muy elemental, son
de uso constante en casi todas las aplicaciones.

\setcounter{section}{0}

\section{Condiciones de finitud relativas. Conjuntos construibles en los preesquemas}
\label{section:IV.1-fr}

Vamos a retomar, en este parágrafo, completándola, la exposición de las
«condiciones de finitud» para un morfismo de preesquemas $f:X\to Y$, dada
en (\textbf{I}, 6.3 y 6.6). Hay esencialmente dos nociones de «finitud» de
naturaleza \emph{global} sobre $X$, la de morfismo \emph{cuasicompacto}
(definida en (\textbf{I}, 6.6.1)) y la de morfismo \emph{cuasiseparado}; hay,
por otra parte, dos nociones de «finitud» de naturaleza \emph{local} sobre $X$,
la de morfismo \emph{localmente de tipo finito} (definida en (\textbf{I}, 6.6.2))
y la de morfismo \emph{localmente de presentación finita}. Combinando
estas nociones locales con las nociones globales anteriores, se obtienen las nociones
de morfismo \emph{de tipo finito} (definida en (\textbf{I}, 6.3.1)) y de morfismo
\emph{de presentación finita}. Para comodidad del lector, daremos de
nuevo, en este parágrafo, las principales propiedades enunciadas en (\textbf{I}, 6.3
y 6.6), remitiendo, por supuesto, a esos números del capítulo I para sus
demostraciones.

En los n\textsuperscript{os}~(1.8) y (1.9), completamos, en el marco de los
preesquemas, y haciendo uso de las nociones de finitud anteriores, los
resultados sobre los conjuntos construibles dados en
(\textbf{0}\textsubscript{III}, §~9).

